% Диалектический характер цикла
\section{Dialectical nature of the cycle}

\begin{frame}
  \frametitle{Mechanism of dialectical development}
  The mechanism of the dialectical development of the system includes four successive stages:
  \begin{enumerate}
    \item \textit{Accumulation} of contradictions of the system in a hidden form through gradual changes in the quantitative ratio of its dialectical opposites;
    \item \textit{Aggravation} of the accumulated contradictions of the system in an explicit form to a critical level due to a significant change in the quantitative ratio of its dialectical opposites;
    \item \textit{Resolution} of contradictions incompatible with the existence of the current system by violent means by breaking the dialectical unity of its opposites during the crisis;
    \item \textit{Removal} of contradictions of the previous system in the course of its denial or transition to a qualitatively new level of development through the formation of a new dialectical unity of opposites.
  \end{enumerate}

  \begin{alertblock}{Finding}
    \textbf{Crisis} – natural turning point, acting as a constitutive phase of the cyclical development of a system, expressed in the rupture of the former form of unity of its opposites and thereby causing its transition to a qualitatively new state within each subsequent turn of the spiral of dialectical development.
  \end{alertblock}
\end{frame}

\begin{frame}
  \frametitle{Cycle as a form of dialectical development}
  The fundamental contradictions of all constituent elements of the capitalist economy within its framework can be classified into two main categories:
  \begin{itemize}
    \item \textit{Main contradictions} – directly define the capitalist economy, cannot be resolved within the framework of capitalism and form the basis of more specific secondary contradictions;\footnote{In particular, the contradiction between \textit{social character of the productive forces} and \textit{private ownership of the means of production}}
    \item \textit{Secondary contradictions} – lie on the surface, representing a more concrete and derivative form of the main contradictions, partially and temporarily resolving by force during the crisis of the capitalist economy, reproducing at a higher stage of its development.\footnote{In particular, the contradiction between \textit{high production potential of capitalism} and \textit{low consumption capacity of the masses}}
  \end{itemize}

  \begin{alertblock}{Finding}
    The movement of the capitalist economy should be characterized precisely as an \textbf{industrial cycle}, since the temporary partial resolution of its \textit{secondary contradictions} during the crisis does not entail the negation of the capitalist system itself as such, while the \textit{main contradictions} serving as the primary source of its movement remain unchanged at all stages of its development.
  \end{alertblock}
\end{frame}